\documentclass[11pt]{article}
\usepackage{amsmath,amssymb,amsthm}
\usepackage[margin=2.5cm]{geometry}
\newtheorem{theorem}{Theorem}
\newtheorem{proposition}[theorem]{Proposition}
\newtheorem{lemma}[theorem]{Lemma}
\theoremstyle{remark}
\newtheorem{remark}[theorem]{Remark}
\title{C\_hunt: proofs behind the hunt}
\date{}
\begin{document}
\maketitle

\noindent Notation is that of \texttt{unify/BRIEF.md}.  $p$ is an odd prime, $\rho=e^{2\pi i/p}$, $F=\{1,\dots,(p-1)/2\}$, $n\ge2$,
$\nu=n+1$, $b\in[2,\nu]$, $s=\nu/b$, $r=1/s$, $\kappa_l(\delta)=1-\cos(2\pi l\delta/p)$ and $c_l(\delta)=1-\rho^{-l\delta}$, so that
$|c_l(\delta)|^2=2\kappa_l(\delta)$ and $c_l(-\delta)=\overline{c_l(\delta)}$.

\section{The class-$\mathcal C$ Riesz-product model}
A class-$\mathcal C$ neuron has a single frequency $l\in F$ in every slot and no input DC.  Output DC cancels in every gap and only adds
cost, so we may drop it.  Balancing the slot norms keeps the table and does not increase the cost, so for upper bounds
we may take every neuron balanced.  Write the neuron as $u_t(a)=\alpha\cos(2\pi la/p+\varphi_t)$, $w(c)=\alpha\cos(2\pi lc/p+\psi)$.
A direct computation from $\Phi_l(x)=\frac2p\rho^{l y(x)}\sum_cf(x)[c]\rho^{-lc}$ gives
\[
\Phi_l(x)=b\prod_{t=1}^n\big(1+\xi_t\zeta_l^{x_t}\big),\qquad \zeta_l=\rho^{2l},\ \ \xi_t=e^{2i\varphi_t},\ \ b=2^{-n}\alpha^{\nu}e^{i(\psi-\sum_t\varphi_t)},
\]
and $|b|=B_p\|\omega\|^{\nu}$.  Every $(b,\xi)\in\mathbb C\times\mathbb T^n$ occurs.  Summing over the neurons $i$ of frequency $l$:
\[
\Phi_l=A_l+f_l,\qquad A_l=\sum_ib_i,\qquad f_l=\sum_{\emptyset\ne T\subseteq[n]}a_{l,T}\,\chi_{l,T},\qquad a_{l,T}=\sum_ib_i\prod_{t\in T}\xi_{it},
\qquad \chi_{l,T}(x)=\zeta_l^{\sum_{t\in T}x_t}.
\]
The gaps are $g(x,\delta)=\sum_l\Re\big(\Phi_l(x)c_l(\delta)\big)$.  In units $B_p=1$ the cost is $P=\sum_l Q_l$ with $Q_l=\sum_{i\in l}|b_i|^r$,
and $\gamma=B_p\,h/P^s$.  We also write
\[
G_l=\frac{|A_l|}{\nu^{1-s}Q_l^s}\ \ \text{(gain)},\qquad \varepsilon_l=\frac{\sum_T|a_{l,T}|}{|A_l|}\ \ \text{(relative $L^1$ fluctuation)}.
\]
\emph{Check.} \texttt{t0\_lib.py}: model gaps against the full-table gaps of the refereed GLOBAL library agree to $10^{-14}$
for $(p,n)\in\{(5,2),(7,2),(11,2),(5,3),(7,3),(5,4)\}$; the clean coefficients agree to $4\cdot10^{-15}$; the costs agree to $4\cdot10^{-16}$.

\section{The $L^1$ reduction}
Fix $\theta\ge0$ and consider two statements.

\medskip\noindent$(J'_\theta)$ at an offset $\delta$:\quad
$\displaystyle \min_{x\in\mathbb Z_p^n,\ \sigma=\pm1} g(x,\sigma\delta)\ \le\ \sum_l|A_l|\kappa_l(\delta)\ -\ \theta\sum_{l}\sum_{T}|a_{l,T}|\,|c_l(\delta)|.$

\smallskip\noindent$(P_\theta)$ for a feature $l$ ($A_l\neq0$):\quad $G_l\,(1-\theta\varepsilon_l)\le1$.

\medskip\noindent
$(P_\theta)$ concerns one feature's atom system and does not involve $p$.  $(J'_\theta)$ is a statement about one trigonometric
polynomial on $\mathbb Z_p^n$, evaluated at the pair $\pm\delta$.  It charges the fluctuations of every frequency, ghosts included
(a ghost is a frequency with neurons but $A_l=0$).

\begin{theorem}[$L^1$ reduction; every $p$, $n$, $b$, $S$]\label{thm:red}
Let $\theta$ be a balanced class-$\mathcal C$ network whose features ($A_l\ne0$) lie in $S\subseteq F$.  Put $Q_{\rm tot}=\sum_{l\in F}Q_l$,
$y_l=(Q_l/Q_{\rm tot})^s$ for $l\in S$, and let $\delta^\ast$ minimise $\sum_{l\in S}y_l\kappa_l(\delta)$ over $\delta\ne0$.
If every feature satisfies $(P_\theta)$ and the network satisfies $(J'_\theta)$ at $\delta^\ast$, then
\[\gamma(\theta)\ \le\ \nu^{1-s}B_p\,\Gamma^{(S)}_s(p)\ \ \le\ D_S=\pi_{n,b}(p)B_p\Gamma^{(S)}_s(p).\]
Consequently, if $\theta^\ast(p,n):=\inf(\text{best constant in }(J'))\ \ge\ \theta_P(n,s):=\sup_{G>1}(1-1/G)/\varepsilon$, then
$\mathrm{RD}_S$, and indeed the class-$\mathcal C$ interference-free value on $S$, holds at $(p,n,s)$ for every $S$.
\end{theorem}
\begin{proof}
$h\le\min_{x,\sigma}g(x,\sigma\delta^\ast)$ because $h$ is the minimum over all $(x,\delta)$.  By $(J'_\theta)$,
\[h\le\sum_l\Big(|A_l|\kappa_l(\delta^\ast)-\theta\sum_T|a_{l,T}|\,|c_l(\delta^\ast)|\Big).\]
For a ghost, the summand is $\le0$.  For a feature, $|c_l|=\sqrt{2\kappa_l}\ge\kappa_l$ because $\kappa_l\le2$, so the summand is at most
$|A_l|\kappa_l(1-\theta\varepsilon_l)\le|A_l|\kappa_l(1-\theta\varepsilon_l)_+$.
If $1-\theta\varepsilon_l>0$, then $(P_\theta)$ gives $|A_l|(1-\theta\varepsilon_l)=\nu^{1-s}Q_l^sG_l(1-\theta\varepsilon_l)\le\nu^{1-s}Q_l^s$.
Hence $h\le\nu^{1-s}\sum_{l\in S}Q_l^s\kappa_l(\delta^\ast)=\nu^{1-s}Q_{\rm tot}^s\sum_{l\in S}y_l\kappa_l(\delta^\ast)$.
Since $\sum_{l\in S}y_l^r=\sum_{l\in S}Q_l/Q_{\rm tot}\le1$, the vector $y$ is admissible for the programme on $S$ (with budget at most $1$), so
$\min_\delta\sum_Sy_l\kappa_l(\delta)\le\Gamma^{(S)}_s$.  Finally $\gamma=B_ph/Q_{\rm tot}^s$, and $\pi\ge\nu^{1-s}$.
\end{proof}
\emph{Check.} \texttt{check\_reduction.py}, output \texttt{out/check\_reduction.out}.  It tests $1073$ class-$\mathcal C$ networks: random, IF-perturbed,
cheap-profile, partially optimised, and the hunt's saved optima, at $(p,n,s)\in\{(5,2,\frac32),(7,2,\frac32),(11,2,\frac32),(13,2,1.2),(7,3,2),(5,3,\frac43),(5,4,\frac52)\}$
and at $p\le31$.  For each it computes the network's own constants
$\theta_{\rm net}$ (the best $\theta$ in $(J')$ at $\delta^\ast$) and $\theta_{\rm pay}=\max_l(1-1/G_l)/\varepsilon_l$.  Every
inequality of the chain is checked separately.  Result: failures $=0$.

\begin{remark}
The only $p$-dependent input is $(J')$, at a single pair of offsets.  The only atom input is $(P)$, which is free of $p$.
Neither uses the size of $p$ or of $n$; nothing is sent to a limit.  The conclusion is the interference-free value
\emph{within class $\mathcal C$}.  This is consistent with $p=3$: the certified $p=3$ winners beat the IF value only through DC input
components, which class $\mathcal C$ excludes.  So the theorem asserts nothing about them.
\end{remark}

\section{Single-character cheap features at $n=2$}
\begin{proposition}\label{prop:single}
Let $n=2$ and $0<c<2$.  Take three atoms with masses $|b_i|$ proportional to $(1,1,c)$, aligned ($b_i>0$), with slot-2 phases
$\xi_{i2}=(e^{i\beta},e^{-i\beta},-1)$ where $2\cos\beta=c$, and slot-1 phases $\xi_{i1}=e^{i\phi}\,\overline{\xi_{i2}}^{\,2}$.  Then
\[a_{(0,1)}=0,\qquad a_{(1,1)}=0,\qquad a_{(1,0)}=-(1-c)\,A\,e^{i\phi},\]
so the fluctuation is a single character of relative size exactly $1-c$.  With $\xi_{i1}=e^{i\phi}\xi_{i2}$ instead, the
fluctuation is $-(1-c)Ae^{i\phi}\chi_{(1,1)}$.  The gain is
\[G(c)=3^{s-1}\frac{2+c}{(2+c^{r})^{s}},\qquad G(0^+)=(3/2)^{s-1},\qquad \lim_{c\to0}\frac{1-1/G(c)}{1-c}=1-(2/3)^{s-1}=c_{12}.\]
Hence $\theta_P(2,s)\ge c_{12}$.  The value $c_{12}$ is attained by two aligned atoms (the case $c=0$: masses $(\frac12,\frac12)$, $\xi_{\cdot2}=(1,-1)$, $\xi_{\cdot1}=(1,1)$).
\end{proposition}
\begin{proof}
The masses $(1,1,c)$ close a triangle with the unit vectors $(e^{i\beta},e^{-i\beta},-1)$, because $2\cos\beta-c=0$; this gives $a_{(0,1)}=0$.
Next, $a_{(1,1)}=e^{i\phi}\sum_ib_i\overline{\xi_{i2}}=e^{i\phi}\overline{a_{(0,1)}}=0$.  And
$a_{(1,0)}=e^{i\phi}\sum_ib_i\overline{\xi_{i2}}^{\,2}\propto e^{i\phi}(e^{-2i\beta}+e^{2i\beta}+c)=e^{i\phi}(c^2+c-2)$, while $A\propto2+c$.
Since $c^2+c-2=-(1-c)(2+c)$, this is the claim.  The second variant is identical, with $\sum_ib_i\xi_{i2}^2$.
For the gain, write $w=(1,1,c)/(2+c)$: then $Q=\sum w_i^r=(2+c^r)/(2+c)^r$ and $G=1/(3^{1-s}Q^s)$.  At $c=0$, $G=(3/2)^{s-1}$ and $\varepsilon=1$.
\end{proof}
\emph{Check.} \texttt{hunt\_fejer.py} prints the realised relative fluctuations $0.600000$, $10^{-16}$, $10^{-16}$ for $c=0.4$ at $p=13$.
\texttt{ptest.py} (adversarial maximisation of $(1-1/G)/\varepsilon$ over atom systems with $2\le m\le2\nu+1$ atoms) returns $\theta_P(2,\frac32)=0.183503$, which is $c_{12}$ to $10^{-6}$,
attained at $m=2$ with masses $(\frac12,\frac12)$.

\begin{proposition}[punctured IF block; every $n$]\label{prop:punct}
Let $n\ge1$, $\zeta=e^{2\pi i/n}$, and take $n=\nu-1$ aligned atoms $b_k=A/n$ with $\xi_{kt}=\zeta^{k}$ for every slot $t$ ($k=0,\dots,n-1$).  Then
$a_T=A\cdot[\,n\mid |T|\,]$, so the only nonzero fluctuation coefficient is $a_{[n]}=A$: the field is $A(1+\chi_{l,[n]})$, a single character
of relative size $\varepsilon=1$.  The gain is $G=(\nu/n)^{s-1}$, hence
\[\theta_P(n,s)\ \ge\ 1-\Big(\frac{n}{n+1}\Big)^{s-1}\qquad(=c_{12}\text{ at }n=2).\]
\end{proposition}
\begin{proof}
$a_T=\frac An\sum_{k=0}^{n-1}\zeta^{k|T|}$, which is $A$ if $n\mid|T|$ and $0$ otherwise; for $\emptyset\neq T\subseteq[n]$ this happens only for $T=[n]$.
$Q=n(A/n)^r$, so $\nu^{1-s}Q^s=\nu^{1-s}n^{s-1}A$ and $G=(\nu/n)^{s-1}$.  Then $(1-1/G)/\varepsilon=1-(n/\nu)^{s-1}$.
\end{proof}
\emph{Check.} \texttt{ptest.py} returns $\theta_P(2,1.2)=0.077892$, $\theta_P(2,1.5)=0.183503$ and $\theta_P(3,1.5)=0.133974$, equal to
$1-(n/\nu)^{s-1}$ to $10^{-6}$; the maximisers are $n$ equal atoms with $\varepsilon=1$.  Values for $n=4,5$ are in the report.

\section{$(J)$ with constant $\frac12$ is false}
\begin{proposition}[REFUTED: paired joint $L^1$ detection with constant $\frac12$]
At $p=5$, $n=2$, $\delta=1$ there are coefficients $a_{l,T}$ ($l\in\{1,2\}$, $T\in\{(1,0),(0,1),(1,1)\}$; real $A_l=0$) with
\[\max_{\sigma=\pm1}\Big(-\min_x D(x,\sigma)\Big)=0.33097\ldots\cdot\sum_{l,T}|a_{l,T}|\,|c_l(1)|\ <\ \tfrac12\sum_{l,T}|a_{l,T}|\,|c_l(1)|.\]
The coefficients are listed in \texttt{out/j\_counterexample\_p5n2.out}, and the minimum is re-evaluated over the whole grid in 40-digit
arithmetic.  The best constant found at $p=5$ is $0.2422$ (\texttt{jtest2.py}, 150 restarts).
\end{proposition}
The numerical values of $\theta^\ast(p,n)$ and $\theta_P(n,s)$ are reported in the hunt report.
\end{document}
